% Generated from analysis/final_rule.py. Predictor labels below quote its docstring/code.
% The Round v7 final rule. Prediction only: no files, fitting, or outcomes.
% 
% Arm                 State status                     Math/code          QA
% ------------------  -------------------------------  -----------------  -----------------
% pruning             seen size, unseen d in [0.6,0.9]  v53 power          v53 median curve
% pruning             new size or stage                v53 median curve   v53 median curve
% per-channel RTN     seen bits, new state             per-bit dev median per-bit dev median
% per-channel RTN     seen state                       v36 regression     v36 regression
% grouped RTN         seen state                       v69 interpolation  per-config median
% grouped RTN         new state                        per-config median  per-config median
% distillation        new source, smaller same-stage   linear (math),     constant; 2Wiki only
%                                                      constant (code)
% 
% All returns are absolute deployed losses in native-token nats. Pruning/RTN
% add a development response to source L0; KD adds it to STUDENT dense loss.
% Group interpolation uses v69's pre-specified boundary rule: adjacent log-coordinate
% pairs, zero floor only outside the group grid, and no bit extrapolation.
% Selection is downstream: feasible nominal storage, argmin absolute loss per
% capability, or max_c(L_c-source_L0c), with v64's candidate-set heuristic.
% 
% ``inputs`` contains ``models`` (development-fitted objects), scalar ``L0``,
% ``N0``, ``D0`` and arm-specific ``d`` / ``bit`` / ``config``. KD instead uses
% ``student`` with N0, D0 and a per-capability ``dense`` mapping. QA KD requires
% ``qa_distribution='2Wiki'``. New-size/stage status takes precedence over a
% previously seen size. Undefined table cells are rejected, never improvised.
\begin{table}[tb]
\normalfont
\centering\small
\setlength{\tabcolsep}{3pt}\renewcommand{\arraystretch}{1.05}
\caption{The table gives the selection predictor for each compression method, capability and source-state status. A prediction adds the listed response, in nats per native token, to a dense anchor: the source loss for pruning and quantization, the initial student's loss for distillation. Each method is fitted on its own development panel, listed below. The rule was pre-specified before confirmation.}
\label{tab:locked_rule}
\begin{tabular*}{\textwidth}{@{\extracolsep{\fill}}>{\raggedright\arraybackslash\hspace{0pt}}p{\dimexpr 0.190000000\textwidth-1.900000000\tabcolsep\relax}>{\raggedright\arraybackslash\hspace{0pt}}p{\dimexpr 0.130000000\textwidth-1.300000000\tabcolsep\relax}>{\raggedright\arraybackslash\hspace{0pt}}p{\dimexpr 0.160000000\textwidth-1.600000000\tabcolsep\relax}>{\raggedright\arraybackslash\hspace{0pt}}p{\dimexpr 0.220000000\textwidth-2.200000000\tabcolsep\relax}>{\raggedright\arraybackslash\hspace{0pt}}p{\dimexpr 0.170000000\textwidth-1.700000000\tabcolsep\relax}>{\raggedright\arraybackslash\hspace{0pt}}p{\dimexpr 0.130000000\textwidth-1.300000000\tabcolsep\relax}@{}}
\toprule
Compression method & Capability & State status & Response predictor \\
\midrule
% delta = float(np.dot(p["beta"], z) * prune.shape(d, p["gamma"]))
Pruning & Mathematics & Seen source state; new density & Pruning power form \\
% delta = float(np.dot(p["beta"], z) * prune.shape(d, p["gamma"]))
Pruning & Code & Seen source state; new density & Pruning power form \\
% delta = prune.linear_curve(f["models"][c]["median_curve"]["anchors"], d)
Pruning & Question answering & Seen source state; new density & {Pruning development median curve} \\
% delta = prune.linear_curve(f["models"][c]["median_curve"]["anchors"], d)
Pruning & Mathematics & {New size or new stage} & {Pruning development median curve} \\
% delta = prune.linear_curve(f["models"][c]["median_curve"]["anchors"], d)
Pruning & Code & {New size or new stage} & {Pruning development median curve} \\
% delta = prune.linear_curve(f["models"][c]["median_curve"]["anchors"], d)
Pruning & Question answering & {New size or new stage} & {Pruning development median curve} \\
\addlinespace
% delta = channel_delta(f["models"][c], inputs)
{Per-channel round-to-nearest quantization} & Mathematics & Seen state & {Per-bit source regression} \\
% delta = channel_delta(f["models"][c], inputs)
{Per-channel round-to-nearest quantization} & Code & Seen state & {Per-bit source regression} \\
% delta = channel_delta(f["models"][c], inputs)
{Per-channel round-to-nearest quantization} & Question answering & Seen state & {Per-bit source regression} \\
% delta = f["median"][c][str(inputs["bit"])]
{Per-channel round-to-nearest quantization} & Mathematics & New state & {Per-bit development median} \\
% delta = f["median"][c][str(inputs["bit"])]
{Per-channel round-to-nearest quantization} & Code & New state & {Per-bit development median} \\
% delta = f["median"][c][str(inputs["bit"])]
{Per-channel round-to-nearest quantization} & Question answering & New state & {Per-bit development median} \\
\addlinespace
% anchors = {k: float(np.dot(beta, z)) for k, beta in
%                        f["models"][c]["same_input_interpolation"]["anchors"].items()}
\bottomrule\end{tabular*}
\end{table}
\begin{table}[tb]\ContinuedFloat
\normalfont
\centering\small
\setlength{\tabcolsep}{3pt}\renewcommand{\arraystretch}{1.05}
\begin{tabular*}{\textwidth}{@{\extracolsep{\fill}}>{\raggedright\arraybackslash\hspace{0pt}}p{\dimexpr 0.190000000\textwidth-1.900000000\tabcolsep\relax}>{\raggedright\arraybackslash\hspace{0pt}}p{\dimexpr 0.130000000\textwidth-1.300000000\tabcolsep\relax}>{\raggedright\arraybackslash\hspace{0pt}}p{\dimexpr 0.160000000\textwidth-1.600000000\tabcolsep\relax}>{\raggedright\arraybackslash\hspace{0pt}}p{\dimexpr 0.220000000\textwidth-2.200000000\tabcolsep\relax}>{\raggedright\arraybackslash\hspace{0pt}}p{\dimexpr 0.170000000\textwidth-1.700000000\tabcolsep\relax}>{\raggedright\arraybackslash\hspace{0pt}}p{\dimexpr 0.130000000\textwidth-1.300000000\tabcolsep\relax}@{}}
\toprule
Compression method & Capability & State status & Response predictor \\
\midrule
{Grouped round-to-nearest quantization} & Mathematics & Seen state & {Piecewise source interpolation} \\
% anchors = {k: float(np.dot(beta, z)) for k, beta in
%                        f["models"][c]["same_input_interpolation"]["anchors"].items()}
{Grouped round-to-nearest quantization} & Code & Seen state & {Piecewise source interpolation} \\
% anchors = f["models"][c]["median"]["anchors"]
{Grouped round-to-nearest quantization} & Question answering & Seen state & {Per-configuration development median} \\
% anchors = f["models"][c]["median"]["anchors"]
{Grouped round-to-nearest quantization} & Mathematics & New state & {Per-configuration development median} \\
% anchors = f["models"][c]["median"]["anchors"]
{Grouped round-to-nearest quantization} & Code & New state & {Per-configuration development median} \\
% anchors = f["models"][c]["median"]["anchors"]
{Grouped round-to-nearest quantization} & Question answering & New state & {Per-configuration development median} \\
\addlinespace
% float(distill._predict(f["linear"][c], [{**s, "L0": anchor}])[0])
Distillation & Mathematics & New source & {Linear source regression} \\
% float(f["constant"][c])
Distillation & Code & New source & Pre-specified constant \\
% float(f["constant"][c])
Distillation & Question answering & New source & Pre-specified constant \\
\addlinespace
% value = float(inputs["L0"])
Dense & Mathematics & Any valid status & 0 (no loss change) \\
% value = float(inputs["L0"])
Dense & Code & Any valid status & 0 (no loss change) \\
% value = float(inputs["L0"])
Dense & Question answering & Any valid status & 0 (no loss change) \\
\bottomrule\end{tabular*}
\caption[]{The table continues the pre-specified selection rules. Absolute losses are in nats per native token.}
\end{table}
\begingroup\small
\noindent\textbf{Development data.} Pruning: pruning development set, 17 Pythia states, 84 density responses per capability ($0.55\le d<1$, excluding 2.8 billion). Channel quantization: development selection pre-specified paired channel rows; 16/16/4/17/17 states at 8/6/5/4/3 bits. Grouped quantization: grouped-quantization interpolation development grid, 160 million/410 million/1.4 billion at 16 thousand/143 thousand, $b\in\{3,4,5\}$, $g\in\{64,128,256\}$. Fixed-recipe students: fixed-recipe student fixed-recipe Pythia students pre-specified in development selection: 160 million/410 million/1.4 billion at 16 thousand/64 thousand/143 thousand, excluding both selectable students (160 million at step 64 thousand, 410 million at step 64 thousand), leaving seven. Recipe: teacher gpt-5.6-luna, full pool 600, two epochs, strict low-rank adaptation, seed 0. Full rosters and literal code are in the selection decomposition artifact.
\par\noindent\textbf{Exact forms.} The pruning development power response is $(\beta_c^\top z)[(1-d)/0.3]^{\gamma_c}$, with $z=(1,\allowbreak z(\log N_0),\allowbreak z(L_{0,c}),\allowbreak z(\log D_0))$ and development standardization. Its median curve linearly interpolates adjacent per-density medians. per-bit source regression is per-bit ordinary least squares on $(1,\allowbreak z(\log(N_0/10^9)),\allowbreak z(L_{0,c}),\allowbreak z(\log(D_0/10^9)))$; the new-state branch takes the median signed response at that bit. grouped-quantization interpolation interpolation uses per-configuration source-regression anchors $\beta_{c,b,g}^\top z$; its median branch uses per-configuration development medians. Both pass the anchors through the same piecewise bilinear interpolation in $(\log_2(2^{b-1}-1),\log_2(g/128))$. Outside the group grid, use the nearest boundary pair and floor only the extrapolated response at zero; no bit extrapolation. pruning development power and grouped-quantization interpolation source-anchor fits use ridge $10^{-3}$ including the intercept. Distillation math uses $\beta_0+\beta_1z(\log N_S)+\beta_2z(L_{S0,c})+\beta_3z(\log D_S)$; code and Question answering use the arithmetic mean development response for each capability, rather than its median.
\par\noindent\textbf{Domain and selection.} New status means a new state, a new size, a new stage, or a new source, taking precedence over seen size. Pruning requires $0.6\le d\le0.9$; channel round-to-nearest quantization requires a development bit. Distillation requires a smaller same-stage student ($N_S<N_0$, $D_S=D_0$), and Question answering requires the primary question-answering distribution. Undefined cells are rejected. Dense accepts every valid status. Selection minimizes predicted absolute loss or $\max_c(\widehat L_c-L_{0,c})$ over storage-feasible candidates; multi has no separately fitted predictor. independent selection invokes the new-stage branch for all four states included in the independent confirmation panel.
\par\noindent\textbf{Differences from the delivered-predictor table.} The delivered-predictor table allows power or per-density regression for seen-size pruning; this rule uses power for math/code and medians for Question answering and new size/stage. The channel-quantization branches now agree: both take the per-bit source regression on states in the fit and the per-bit development median on new ones. This rule was pre-specified with that branch before the confirmation panel was measured, and is unchanged. Grouped branches agree, with the boundary rule specified above. Its distillation exposure and joint budget--pool forms are not used here: selection uses the fixed-recipe student-state math form together with the code and Question answering constants fixed at the same time.
\endgroup
